% !PRINTED 630
\weberpage{630}
%<<CONT
cordingly the functions (12) form a basis of the field $\Omega$, and it follows that
\weq{13}{ef = n.}
Thus $e$ and $f$ are divisors of $n$.

If
\weq{14}{\xi_1, \xi_2, \dots, \xi_e}
denotes a basis of the field $\Omega_1$, then the $n$ quantities:
\weq{15}{\xi_1 \theta^k, \xi_2 \theta^k, \dots, \xi_e \theta^k, \qquad k = 0, 1, 2, \dots, f-1}
form a basis of $\Omega$. If one forms the substitution of the $e$ quantities:
\weq{16}{(\xi \xi_i) = Y_1(\xi_i),}
then
\weq{17}{(-1)^e b_e' = |Y_1| = N_1(\xi)}
is the norm of $\xi$ in the field $\Omega_1$, and if one forms the substitution for $(\xi \xi_i \theta_k)$ in the field $\Omega$ with the basis (15), there results
% [sic] printed (\xi \xi_i \theta_k) with subscript k, a misprint: the basis (15) and the display below have the power \theta^k
\[
(\xi \xi_i \theta^k) = Y_1(\xi_i)\, Y_1(\xi_i \theta) \dots Y_1(\xi_i \theta^{e-1}),
% [sic] printed \theta^{e-1} in the last factor, an error: k runs only to f-1 in (15), so the product has f factors ending in Y_1(\xi_i \theta^{f-1}), as N(\xi) = |Y_1|^f in (18) requires
\]
from which
\weq{18}{N(\xi) = |Y_1|^f.}
From the substitution determinant $|Y_1|$ one derives the function $\varphi_1(t)$ by formula (2), and so it follows from (18):
\weq{19}{N(t - \xi) = [\varphi_1(t)]^f,}
and from this:

6. \gesperrt{The function $\varphi(t) = N(t - \xi)$ is either irreducible or a power of an irreducible function.}

If the exponent $f$ of this power is greater than 1, then the field $\Omega$ is imprimitive, $\Omega_1$ is a subfield, and $\Omega$ is a field of degree $f$ over $\Omega_1$ (Bd. I, \S\ 151).

If $f = 1$, $e = n$, then $\xi$ is called a \gesperrt{primitive function of the field} $\Omega$, and $\Omega_1$ is identical with $\Omega$.

We return to equation (2) or (3), which, as we have seen, is independent of the choice of the basis, and define the coefficient
\weq{20}{-b_1 = y_{1,1} + y_{2,2} + \cdots + y_{n,n} = S(\xi)}
as the \gesperrt{trace of the function} $\xi$. For the trace there result from the definition the following fundamental theorems, in which $\xi, \xi'$ denote any two functions in $\Omega$, and $x$ a rational function:
%CONT>>
